\documentclass[11pt, a4paper]{letter}
\usepackage[utf8]{inputenc}
\usepackage{geometry}
\geometry{margin=1in}

\signature{Rafael D. De Paz \\ Independent Researcher \\ Logos Kernel Architect}
\address{Rafael D. De Paz \\ \texttt{rdepaz.com/research} \\ \texttt{nexus@rddp.ai}}

\begin{document}
\begin{letter}{
The Editorial Board \\
Journal of the ACM (JACM) \\
Association for Computing Machinery \\
New York, NY
}

\date{\today}
\opening{Dear Editors,}

I am pleased to submit the enclosed manuscript, \textit{The Thermodynamic Church-Turing Thesis: An Entropic Proof that $\mathbf{P} \neq \mathbf{NP}$}, for consideration as an original research article in the Journal of the ACM.

This paper establishes the mathematical separation of $\mathbf{P}$ and $\mathbf{NP}$ by anchoring abstract algorithmic complexity theory directly into non-equilibrium statistical mechanics. Instead of relying on structural Turing machine constraints, this paper models the problem space as a non-convex pseudo-Riemannian manifold. By unifying Morse Theory, Kolmogorov complexity, and Landauer’s limit, the proof demonstrates that bridging the gap between polynomial-time verification and non-deterministic discovery requires a massive thermodynamic energy penalty—an exact physical tax that is irreconcilable with $O(n^c)$ polynomial time algorithms.

We respectfully assert that this manuscript resolves the longest-standing open problem in theoretical computer science. Due to the cross-disciplinary nature of the proof (computational complexity, topology, and thermodynamic physics), we hope you will engage reviewers with expertise across both statistical mechanics and abstract algorithmic complexity.

This manuscript is original, has not been published previously, and is not currently under consideration by any other journal. We look forward to your rigorous evaluation.

\closing{Sincerely,}

\end{letter}
\end{document}
